\section{Results and Discussion}
\label{sec:results}

\subsection{Overview}

This section reports the results of the two-phase framework and organizes the
discussion around four research questions:
\begin{itemize}
\item \textbf{RQ1:} Which handcrafted descriptor configurations provide the
highest class separability?
\item \textbf{RQ2:} Which feature representations achieve the best classification
performance?
\item \textbf{RQ3:} Which classifiers are most effective on the optimized
handcrafted features?
\item \textbf{RQ4:} Are the observed performance differences statistically
significant?
\end{itemize}

\subsection{Phase 1: Separability-Guided Feature Configuration}

\subsubsection{Descriptor Configuration Ranking (RQ1)}
For each descriptor family, the candidate parameter combinations were ranked by
the composite separability score of Section~\ref{sec:phase1}. The selected
configurations and their scores are reported in Table~\ref{tab:phase1}. Across
all four families, the scores varied substantially between candidate settings,
confirming that descriptor configuration has a marked effect on feature quality.
LBP, DWT, and HOG each achieved composite scores above $0.98$ within their
respective search spaces, whereas GLCM exhibited greater sensitivity to
parameter choice (best score $0.747$). As noted in Section~\ref{sec:phase1}, the
composite scores are normalized within each descriptor family and therefore rank
configurations within a family rather than across families.

\subsubsection{Analysis of Separability Metrics}
The three metrics offered complementary views of feature quality. The Fisher
Discriminant Ratio favored configurations with strong inter-class separation and
low intra-class variance; Mutual Information favored configurations preserving
the most class-relevant information; and the Davies--Bouldin Index favored
compact, well-separated class clusters. Combining the three yielded a more
balanced ranking than any single criterion.

\subsubsection{Discussion}
The optimization confirmed that the discriminative power of handcrafted
descriptors is sensitive to parameter selection, that different families differ
in this sensitivity, and that the selected configurations formed the basis for
the classification experiments that follow.

\subsection{Phase 2: Classifier Benchmarking}

\subsubsection{Overall Benchmark Performance (RQ3)}
Table~\ref{tab:overall} summarizes the mean cross-validation performance of
each classifier averaged across all eight feature sets. Support Vector Machines
and Extra Trees achieved the highest mean macro-F1 ($0.829$), followed closely
by XGBoost ($0.826$), LightGBM ($0.822$), and Random Forest ($0.822$).
Distance-based KNN and linear logistic regression ranked lower on average
($0.793$ and $0.792$, respectively), although KNN reached near-top performance
on the strongest feature sets (Table~\ref{tab:leaderboard}).

\begin{table}[t]
\centering
\caption{Mean cross-validation performance per classifier, averaged across all
feature sets (5-fold stratified CV).}
\label{tab:overall}
\renewcommand{\arraystretch}{1.2}
\begin{tabular}{@{}lcccc@{}}
\toprule
\textbf{Classifier} & \textbf{Acc.} & \textbf{Prec.} & \textbf{Rec.} & \textbf{F1} \\
\midrule
SVM      & $0.829$ & $0.832$ & $0.829$ & $0.829$ \\
ET       & $0.830$ & $0.833$ & $0.830$ & $0.829$ \\
XGBoost  & $0.827$ & $0.828$ & $0.827$ & $0.826$ \\
LightGBM & $0.822$ & $0.823$ & $0.822$ & $0.822$ \\
RF       & $0.823$ & $0.825$ & $0.823$ & $0.822$ \\
MLP      & $0.810$ & $0.813$ & $0.810$ & $0.810$ \\
KNN      & $0.794$ & $0.798$ & $0.794$ & $0.793$ \\
LR       & $0.793$ & $0.794$ & $0.793$ & $0.792$ \\
\bottomrule
\end{tabular}
\end{table}

\subsubsection{Feature Set Comparison (RQ2)}
Table~\ref{tab:featurescore} reports the mean performance for each feature
representation averaged across all eight classifiers. The full fused
representation $\mathrm{Full(opt)}$ achieved the highest mean macro-F1
($0.887$), followed by the gradient-only set $D$ ($0.865$). Statistical
features alone ($A$) and the frequency-domain set $C$ yielded the lowest mean
scores ($0.765$ and $0.717$, respectively). Intermediate fusion sets
($A{+}B$, $A{+}B{+}C$, $B{+}C$) improved over their weakest constituents but
did not surpass $D$ or $\mathrm{Full(opt)}$, indicating that optimized HOG
descriptors carry the strongest discriminative signal in this benchmark.

\begin{table}[t]
\centering
\caption{Mean cross-validation performance per feature representation, averaged
across all classifiers (5-fold stratified CV).}
\label{tab:featurescore}
\renewcommand{\arraystretch}{1.2}
\begin{tabular}{@{}lcc@{}}
\toprule
\textbf{Feature set} & \textbf{Accuracy} & \textbf{F1} \\
\midrule
Full(opt)  & $0.887$ & $0.887$ \\
D          & $0.865$ & $0.865$ \\
A+B        & $0.835$ & $0.835$ \\
A+B+C      & $0.828$ & $0.828$ \\
B+C        & $0.822$ & $0.822$ \\
B          & $0.805$ & $0.804$ \\
A          & $0.765$ & $0.765$ \\
C          & $0.720$ & $0.717$ \\
\bottomrule
\end{tabular}
\end{table}

\begin{figure}[t]
\centering
\includegraphics[width=\linewidth]{p2_f1_best_per_featureset}
\caption{Best test macro-F1 per feature representation (bar colour indicates
the winning classifier).}
\label{fig:phase2_best_per_fs}
\end{figure}

\begin{figure*}[t]
\centering
\includegraphics[width=\textwidth]{p2_f1_heatmap}
\caption{Phase~2 benchmark heatmap: test macro-F1 for every
classifier--feature-set pair.}
\label{fig:phase2_dashboard}
\end{figure*}

\subsubsection{Leaderboard of Best Combinations}
Table~\ref{tab:leaderboard} ranks the top ten classifier--feature-set
combinations by macro-F1. The best overall configuration pairs a polynomial-kernel
SVM with $\mathrm{Full(opt)}$ (accuracy $=0.920$, macro-F1 $=0.920$,
$\kappa=0.88$). A second SVM on the HOG-only set $D$ ranks second
(F1 $=0.903$), and KNN on $\mathrm{Full(opt)}$ ranks third (F1 $=0.900$).
Notably, six of the top ten entries use $\mathrm{Full(opt)}$, confirming the
benefit of multi-family fusion when paired with an appropriate classifier.

% Auto-generated — do not edit by hand
\begin{table}[!t]
\centering
\caption{Top-10 classifier--feature-set combinations ranked by macro-F1 (5-fold stratified CV). $\kappa$: Cohen's kappa; F1 CI: bootstrap 95\% interval.}
\label{tab:leaderboard}
\renewcommand{\arraystretch}{1.15}
\footnotesize
\setlength{\tabcolsep}{3pt}
\begin{tabular}{@{}rllcccl@{}}
\toprule
\textbf{Rank} & \textbf{Feature set} & \textbf{Clf.} & \textbf{Acc.} & \textbf{F1} & $\boldsymbol{\kappa}$ & \textbf{F1 CI} \\
\midrule
1 & Full(opt) & SVM & 0.9200 & 0.9204 & 0.88 & [0.8895, 0.9504] \\
2 & D --- HOG & SVM & 0.9033 & 0.9032 & 0.855 & [0.8683, 0.9346] \\
3 & Full(opt) & KNN & 0.9000 & 0.9002 & 0.85 & [0.866, 0.9322] \\
4 & D --- HOG & KNN & 0.8967 & 0.8967 & 0.845 & [0.8626, 0.9283] \\
5 & Full(opt) & LightGBM & 0.8933 & 0.8938 & 0.84 & [0.8592, 0.9267] \\
6 & Full(opt) & MLP & 0.8833 & 0.8839 & 0.825 & [0.8457, 0.9187] \\
7 & Full(opt) & XGBoost & 0.8800 & 0.8800 & 0.82 & [0.8427, 0.9145] \\
8 & Full(opt) & ET & 0.8767 & 0.8765 & 0.815 & [0.84, 0.9123] \\
9 & D --- HOG & MLP & 0.8767 & 0.8767 & 0.815 & [0.8382, 0.9119] \\
10 & Full(opt) & LR & 0.8700 & 0.8702 & 0.805 & [0.8306, 0.905] \\
\bottomrule
\end{tabular}
\end{table}


\begin{figure}[t]
\centering
\includegraphics[width=\linewidth]{p2_leaderboard}
\caption{Phase~2 leaderboard: macro-F1 for all classifier--feature-set
combinations.}
\label{fig:phase2_leaderboard}
\end{figure}

\subsubsection{Best Classifier--Feature-Set Combinations}
Table~\ref{tab:best} highlights the three highest-ranked combinations from the
full benchmark (64 configurations).

\begin{table}[t]
\centering
\caption{Top three classifier--feature-set combinations (5-fold stratified CV).}
\label{tab:best}
\renewcommand{\arraystretch}{1.2}
\begin{tabular}{@{}cllcc@{}}
\toprule
\textbf{Rank} & \textbf{Feature set} & \textbf{Classifier} & \textbf{Acc.} & \textbf{F1} \\
\midrule
1 & Full(opt) & SVM      & $0.920$ & $0.920$ \\
2 & D         & SVM      & $0.903$ & $0.903$ \\
3 & Full(opt) & KNN      & $0.900$ & $0.900$ \\
\bottomrule
\end{tabular}
\end{table}

\subsection{Per-Feature-Set Benchmark Results}

Rather than stacking metrics, hyperparameter grids, and confusion matrices into
one panel per feature set (which is unreadable at paper scale), we use three
complementary views: the heatmap in Figure~\ref{fig:phase2_dashboard}, the
best-model-per-set summary in Figure~\ref{fig:phase2_best_per_fs}, and focused
bar charts for the two strongest representations ($\mathrm{Full(opt)}$ and $D$).

\begin{figure}[t]
\centering
\includegraphics[width=\linewidth]{p2_f1_bars_full}
\caption{Test macro-F1 of all eight classifiers on $\mathrm{Full(opt)}$.}
\label{fig:p2_full_bars}
\end{figure}

\begin{figure}[t]
\centering
\includegraphics[width=\linewidth]{p2_f1_bars_hog}
\caption{Test macro-F1 of all eight classifiers on $D$ (HOG only).}
\label{fig:p2_hog_bars}
\end{figure}

\subsection{Class-Wise Performance Analysis}

\subsubsection{Confusion Matrix}
Figure~\ref{fig:cm} shows the row-normalized confusion matrix for the best
overall configuration (SVM on $\mathrm{Full(opt)}$, test split).

\begin{figure}[t]
\centering
\includegraphics[width=0.72\linewidth]{p2_confmat_best}
\caption{Confusion matrix for SVM on $\mathrm{Full(opt)}$ (test set).}
\label{fig:cm}
\end{figure}

\subsubsection{Per-Class Metrics (RQ2)}
Table~\ref{tab:perclass} reports aggregate precision, recall, and macro-F1 for
the best combination across all three classes. Per-class breakdown from the
confusion matrix will be added once the figure is inserted.

\begin{table}[t]
\centering
\caption{Aggregate performance for the best combination (SVM on Full(opt)),
5-fold CV.}
\label{tab:perclass}
\renewcommand{\arraystretch}{1.2}
\begin{tabular}{@{}lccc@{}}
\toprule
\textbf{Metric} & \textbf{Precision} & \textbf{Recall} & \textbf{F1} \\
\midrule
Macro (all classes) & $0.923$ & $0.920$ & $0.920$ \\
\bottomrule
\end{tabular}
\end{table}

\subsection{Feature Contribution Analysis}

The feature-set comparison (Table~\ref{tab:featurescore}) indicates that
gradient-based HOG descriptors dominate discriminative performance: $D$ alone
achieves $0.865$ mean macro-F1, substantially outperforming texture ($B$,
$0.804$) and frequency ($C$, $0.717$) families. Adding statistical, texture,
and frequency cues to HOG in $\mathrm{Full(opt)}$ yields a further gain to
$0.887$ mean macro-F1 and the single best configuration (SVM, F1 $=0.920$).
Tree-ensemble feature-importance analysis can further quantify branch-level
contributions; this is left for future interpretability reporting.

\subsection{Relationship Between Separability and Classification Performance}

Phase~1 composite scores are normalized within each descriptor family and are
not directly comparable across families. Qualitatively, families with high
within-family separability scores (LBP, DWT, HOG above $0.98$) correspond to
stronger Phase~2 performance when used in fusion or alone ($B$, $D$,
$\mathrm{Full(opt)}$), whereas GLCM's lower separability score ($0.747$) aligns
with texture-only representations ranking below gradient-based sets. A formal
correlation analysis between raw per-configuration separability metrics and
downstream classification F1 remains a direction for future work.

\subsection{Comparison with Existing Approaches}

The best cross-validated configuration (SVM on $\mathrm{Full(opt)}$, macro-F1
$=0.920$) provides a strong handcrafted-feature baseline on the three-class
Cheng et al.\ dataset~\cite{cheng2017}. Comparisons with prior studies should
account for differences in class count, evaluation protocol, and whether
results are reported on cross-validation or a held-out split
(e.g.,~\cite{pattanaik2022,dheepak2024,nawaz2022,basthikodi2024}). The
methodological contributions that distinguish this framework are its systematic,
classifier-independent descriptor optimization; multi-family feature fusion;
comprehensive feature cleaning; broad eight-classifier benchmark; and rigorous
statistical validation (Section~\ref{sec:statistics}).

\subsection{Practical Implications}
\label{sec:practical}

\subsubsection{Small-Scale and Limited-Data Settings}
Many medical imaging datasets remain relatively small due to annotation costs and
privacy constraints. Handcrafted feature engineering remains particularly attractive
under such conditions because it can achieve strong performance without requiring
large-scale training data.

\subsubsection{Reproducibility}
The explicit feature extraction and optimization procedures contribute to
reproducible experimentation and facilitate independent validation by future
researchers.

\subsection{Discussion}

Three principal conclusions emerge from the benchmark. First, descriptor
parameter optimization has a substantial impact on feature discriminability and
classification performance; feature extraction should be treated as an integral
component of the learning pipeline rather than a fixed preprocessing step.
Second, optimized HOG gradient features provide the strongest single-family
representation, while multi-family fusion in $\mathrm{Full(opt)}$ yields the
best overall results. Third, margin-based and ensemble classifiers (SVM,
Extra Trees, boosting methods) perform most consistently across feature sets,
although instance-based KNN achieves competitive peak performance on the richest
representations.

\subsection{Summary}
\label{sec:results_summary}

The findings presented in this section demonstrate that separability-guided
handcrafted feature optimization constitutes an effective strategy for improving
brain tumor MRI classification. The results highlight the importance of descriptor
configuration, confirm the benefits of feature fusion, and demonstrate that
carefully optimized handcrafted features can support highly competitive machine
learning performance (best macro-F1 $=0.920$). Moreover, the comprehensive
benchmarking and statistical validation framework provides strong evidence that
the observed improvements are meaningful, reproducible, and relevant for practical
computer-aided diagnosis applications. These findings form the basis for the
conclusions presented in Section~\ref{sec:conclusion}.
